% Standalone addition, following the ordered-vector parity proposition.
\begin{proposition}[Proved: an integral free-valued diagram with torsion colimit]
\label{A:ordered-vector-integral-torsion}
There is a diagram $M:\mathcal V^{<}_2\to\Ab$ whose values are
finite-rank free abelian groups and whose colimit contains an element
of order exactly two. This assertion remains valid whether the empty
punctured zero-dimensional vector space is included or omitted.
\end{proposition}
\begin{proof}
For an object $B$, let $\mathcal T(B)$ be the finite set of natural
binary orders on its underlying vector space, with zero omitted.
An embedding $a:B\to C$ induces the restriction map
$r_a:\mathcal T(C)\to\mathcal T(B)$. Restricting a natural order to
a subspace again gives a natural order: an adapted echelon basis
orders the subspace by its successive leading coordinates. Define
\[
 N(B)=\mathbb Z^{\mathcal T(B)},\qquad
 N(a)(f)=f\circ r_a,\qquad
 M(B)=N(B)/\mathbb Z\mathbf1_B.
\]
These are covariant functors. Each $N(B)$ is free of finite rank,
and its constant function $\mathbf1_B$ is primitive, so $M(B)$ is
also free of finite rank.

The uniform probabilities on $\mathcal T(B)$ are restriction
consistent. Indeed every natural order on a subspace extends by
extending an adapted ordered basis. The linear automorphisms of the
ambient space stabilizing that subspace act transitively on its
natural orders, and restriction is equivariant for this action.
Consequently all restriction fibers have the same cardinality.
It follows that
\[
 \mu_B(f)=\frac1{|\mathcal T(B)|}
                \sum_{\tau\in\mathcal T(B)}f(\tau)
\]
defines a natural homomorphism $N\to\mathbb Q$, where the latter
is the constant diagram. Since $\mathcal V^{<}_2$ is connected,
it induces
\[
 \mu:\operatorname{colim}N\longrightarrow\mathbb Q,
 \qquad \mu([\mathbf1])=1.
\]
In particular the map $\mathbb Z\to\operatorname{colim}N$ induced
by constant functions is injective. Right exactness of colimits gives
\[
 \operatorname{colim}M
   =\bigl(\operatorname{colim}N\bigr)/\mathbb Z[\mathbf1].
 \tag{IT.1}
\]

Take the objects $P,Q$ and the five embeddings $a_j:P\to Q$ from
Proposition~\ref{A:ordered-vector-parity}. Regard its five indicator
functions $f_j$ as integer-valued functions of the auxiliary natural
order on $P$. The parity identities proved there say precisely that
\[
 \sum_{j=1}^{5}f_j=\mathbf1_P+2h_P,
 \qquad
 \sum_{j=1}^{5}N(a_j)(f_j)=2h_Q
 \tag{IT.2}
\]
for integer-valued functions $h_P\in N(P)$ and $h_Q\in N(Q)$.
Explicitly $h_P=(\sum_jf_j-1)/2$ and
$h_Q=(\sum_jN(a_j)(f_j))/2$; their integrality follows from those
identities at every natural order, rather than from an averaging
argument. Passing to the colimit in (IT.2) yields
\[
 [\mathbf1]=2v,
 \qquad v=[h_Q]-[h_P]\in\operatorname{colim}N.
\]
Applying $\mu$ gives $\mu(v)=1/2$. Thus $v$ cannot be an integer
multiple of $[\mathbf1]$. Its image in (IT.1) is nonzero and is
killed by two, so has order exactly two, as claimed. The attached
finite certificate verifies (IT.2) on the six natural orders of $P$
and the $168$ natural orders of $Q$, and gives
$\mu_P(h_P)=2/3$ and $\mu_Q(h_Q)=7/6$.
\end{proof}

\paragraph{Scope.}
This is torsion in the ordinary colimit of a pointwise free integral
diagram. It does not assert that $T_1(M)$ has torsion, or determine
any positive derived colimit of this diagram. Although the maps of
$N$ have nonnegative entries in the bases of point indicators,
passing to $M=N/\mathbb Z\mathbf1$ does not supply compatible
nonnegative matrices or lifts to homomorphisms of free monoids.
